% Options for packages loaded elsewhere
\PassOptionsToPackage{unicode}{hyperref}
\PassOptionsToPackage{hyphens}{url}
%
\documentclass[
]{article}
\usepackage{amsmath,amssymb}
\usepackage{lmodern}
\usepackage{iftex}
\ifPDFTeX
  \usepackage[T1]{fontenc}
  \usepackage[utf8]{inputenc}
  \usepackage{textcomp} % provide euro and other symbols
\else % if luatex or xetex
  \usepackage{unicode-math}
  \defaultfontfeatures{Scale=MatchLowercase}
  \defaultfontfeatures[\rmfamily]{Ligatures=TeX,Scale=1}
\fi
% Use upquote if available, for straight quotes in verbatim environments
\IfFileExists{upquote.sty}{\usepackage{upquote}}{}
\IfFileExists{microtype.sty}{% use microtype if available
  \usepackage[]{microtype}
  \UseMicrotypeSet[protrusion]{basicmath} % disable protrusion for tt fonts
}{}
\makeatletter
\@ifundefined{KOMAClassName}{% if non-KOMA class
  \IfFileExists{parskip.sty}{%
    \usepackage{parskip}
  }{% else
    \setlength{\parindent}{0pt}
    \setlength{\parskip}{6pt plus 2pt minus 1pt}}
}{% if KOMA class
  \KOMAoptions{parskip=half}}
\makeatother
\usepackage{xcolor}
\setlength{\emergencystretch}{3em} % prevent overfull lines
\providecommand{\tightlist}{%
  \setlength{\itemsep}{0pt}\setlength{\parskip}{0pt}}
\setcounter{secnumdepth}{-\maxdimen} % remove section numbering
\ifLuaTeX
  \usepackage{selnolig}  % disable illegal ligatures
\fi
\IfFileExists{bookmark.sty}{\usepackage{bookmark}}{\usepackage{hyperref}}
\IfFileExists{xurl.sty}{\usepackage{xurl}}{} % add URL line breaks if available
\urlstyle{same} % disable monospaced font for URLs
\hypersetup{
  hidelinks,
  pdfcreator={LaTeX via pandoc}}

\author{Dietmar Gerald Schrausser}
\date{09.05.2026}


\begin{document}

\hypertarget{foliumblat-iiiv}{%
\section{Foliu{[}m{]}/Blat IIIv}\label{foliumblat-iiiv}}

A \emph{contemporary} English translation is presented from comparing
the existing Latin, Early New High German (ENHG) and English
(translation) texts for better comprehensibility. Especially since the
ENHG text is difficult to understand compared to modern High German, as
a \emph{profound} knowledge of German (as a mother tongue) as well as
already aknowledge of various \emph{dialects} is required.

\hypertarget{of-the-work-of-the-third-day}{%
\subsection{Of the Work of the Third
Day}\label{of-the-work-of-the-third-day}}

\begin{quote}
\textbf{T}Ercio die deus aquas sub firmamento in locum vnum
co{[}n{]}gregauit {[}et{]} appareat arida: vocauitq{[}ue{]} aridam
terram. Congretat{[}i{]}ones vero aquarum maria appellauit. Videns
q{[}uia{]} bonum esset: ait germinet terra herbam virentem: {[}et{]}
semen facientem: {[}et{]} lignum pomiferum facientem fructum iuxta
gen{[}us{]} suum.
\end{quote}

After discussing the firmament's integrity, position and order of the
elements, he briefly reminds us of the confluence of waters into
\emph{one} place and the laws prescribed for the sea, so that it does
not flood the land.\footnote{from the `Heptaplus' (Mirandola \&
  Mirandola, \href{https://doi.org/10.3931/e-rara-61217}{1601}, Book 3).}

If the earth was once invisible and is to come into
\emph{sight}\footnote{`conspectum' and `gesycht'.}, it is necessary that
the waters which are under the heaven or under the middle layer of the
air are \emph{brought together} in one place, into a common mixing
basin, where they flow together according to certain laws, just as it is
practically prescribed for the coasts to join together at one point to
form a \emph{dam}\footnote{`matricem'.}.\footnote{from the `Heptaplus'
  (Mirandola \& Mirandola,
  \href{https://doi.org/10.3931/e-rara-61217}{1601}, Book 2, Chapter 5 /
  4).}

It is untrue that water is not found anywhere in remote and isolated
places, since the \emph{Indian} Sea is separated from the Hircanian Sea,
the \emph{Hircanian} Sea from the Adriatic Sea, the \emph{Adriatic} Sea
from the \emph{Euxiniian} Sea, along with countless remote flowing
\textbf{\emph{springs}} \textbf{\emph{and}} \textbf{\emph{lakes}}, by
weirs\footnote{`intercapedine' and `verre'.}. Since this is a separate
and subdivided collection of sea or river waters, \emph{they should
surely be gathered in one common place}\footnote{`Sed \emph{ideo} ad
  locum vnum congregate aque dicu{[}n{]}tur' und `Aber \emph{darumb}
  werden die wasser an ainige stat versamelt genant'.}.\footnote{from
  the `Heptaplus' (Mirandola \& Mirandola,
  \href{https://doi.org/10.3931/e-rara-61217}{1601}, Book 2, Chapter 2).}
\emph{All}, as \textbf{Solomon} says, tend towards the primary sea, are
led to one place of the oceans and united.\footnote{Referring to King
  Solomon's \emph{most significant} verse in Ecclesiastes 1:7, `Omnia
  flumina intrant in mare, et mare non redundat; ad locum unde exeunt
  flumina revertuntur' and further `ut iterum fluant' (c.f. Jiménez de
  Cisneros, \href{https://doi.org/10.3931/e-rara-46695}{1517}). `All
  rivers flow into the sea, yet the sea is never full; to the place from
  which the rivers come, they return - to flow again {[}and again{]}'.
  \emph{This} describes the endless, cyclical nature of the world.}

When the earth\footnote{`terra' and `erde'.} is then overwhelmed by the
waves of the sea, it is initially not useful or practical for us, but
becomes suitable for animals and also for our purposes when the sea
\emph{recedes}, then its fertility and bountifulness become increasingly
apparent. This is also strongly illustrated by \textbf{Moses} when he
presents it as the \emph{parent}\footnote{`parentem' and `gepererein'.}
of herbs, fruits and trees and describes them as immediately green and
blooming after the accumulation of water.\footnote{German Renaissance
  botanists, such as Otto Brunfels
  (\href{https://www.digitale-sammlungen.de/de/view/bsb11200240?page=1}{1532})
  in his groundbreaking work `Herbarum vivae eicones', frequently quote
  this passage from the `Heptaplus' (Mirandola \& Mirandola,
  \href{https://doi.org/10.3931/e-rara-61217}{1601}, Book 2, Chapter 4)
  to provide theological and philosophical support for their plant
  studies.}

So he placed it (the \emph{Earth}) in the middle of the \emph{World}, as
the \emph{center}, endowed with veins\footnote{`venis'.} of metals such
as gold, silver, bronze\footnote{`ere', aeris.}, copper, tin, lead and
(surpassing all) iron, immediately covered with all kinds of herbs, in
their green ripeness, with herb seeds and trees with the sweetest
fruits. It is also said that on that day he created the most fertile and
glorious \emph{Source}\footnote{`ortum' and `garten'.}, Paradise, with
all kinds of woods and trees, composed of all the delights of springs,
with green plains of trees bearing the most abundant fruit.\footnote{an
  excerpt from the `Supplementum Chronicarum' (s. Foresti,
  \href{https://books.google.com/books?id=ei9TruMbYCkC\&printsec=frontcover}{1492}).}

\hypertarget{references}{%
\subsection{References}\label{references}}

Brunfels, O. (1532). \emph{Herbarum Vivae Eicones: Ad Naturae
Imitationem Summa Cum Diligentia \& Artificio Effigiatae, Una Cum
Effectibus Earundem \ldots{} ; Quibus Adiecta Ad Calcem, Appendix
Isagogica de Usu \& Administratione Simplicium}. Argentorati: apud
Ioannem Schottum.
\url{https://www.digitale-sammlungen.de/de/view/bsb11200240?page=1}

Foresti, G. F. (1492). \emph{Supplementum Chronicarum}. Novariensis:
Bernardinus Rizus.
\url{https://books.google.com/books?id=ei9TruMbYCkC\&printsec=frontcover}

Jiménez de Cisneros, F. (1517). \emph{Biblia Polyglotta Complutensis}.
Complutum: Arnaldo Guillén de Brocar.
\url{https://doi.org/10.3931/e-rara-46695}

Mirandola, G., \& Mirandola, G. F. (1601). \emph{Ioannis Pici,
Mirandulae \ldots{} opera quae extant omnia \ldots{}} Basileae: per
Sebastianum Henricpetri. \url{https://doi.org/10.3931/e-rara-61217}

\end{document}
